From Proposition \ref{prop:premium}, the non-disclosure premium exists if and only if $\mathbb{E}[\theta \mid d=0] > \theta$.

For a high-quality seller ($\theta = 1$), this condition requires $\mathbb{E}[\theta \mid d=0] > 1$. But $\theta \in \{0, 1\}$, so $\mathbb{E}[\theta \mid d=0] \leq 1$ always. Therefore the condition never holds: high-quality sellers never earn a premium from hiding.

The inequality is strict whenever the non-disclosing pool contains any low-quality sellers. In that case $\mathbb{E}[\theta \mid d=0] < 1$, so:
\begin{equation}
    v_1(d=1, \theta=1) - v_1(d=0) = \beta_1 \cdot (1 - \mathbb{E}[\theta \mid d=0]) > 0
\end{equation}

Disclosure yields a strictly higher valuation, and therefore strictly higher profit. Adding hiding costs $c(\theta=1) = c_0 + c_1 \geq 0$ only reinforces this: disclosure dominates hiding.
